\subsection{Physicochemical descriptors: definitions and derivations}
\label{sec:physchem-math}

The discovery audit and the tree-based baseline use four classical
descriptors. We state them precisely and derive the properties we rely on.

\begin{definition}[Net charge at pH $p$]
For residue $a$ with side-chain $\mathrm{p}K_a$, the Henderson--Hasselbalch
equation gives the protonated fraction of a basic group and the
deprotonated fraction of an acidic group as
\begin{equation}
f^{+}_a(p) = \frac{10^{\mathrm{p}K_a}}{10^{p} + 10^{\mathrm{p}K_a}}
\qquad
f^{-}_a(p) = \frac{10^{p}}{10^{p} + 10^{\mathrm{p}K_a}},
\label{eq:hh}
\end{equation}
and the expected net charge of sequence $s = s_1\cdots s_n$ at pH $p$ is
\begin{equation}
Q(s; p) = f^{+}_{\mathrm{N\text{-}term}}(p)
  + \!\!\sum_{a \in \{K,R,H\}}\!\! n_a(s)\, f^{+}_a(p)
  - f^{-}_{\mathrm{C\text{-}term}}(p)
  - \!\!\sum_{a \in \{D,E,C,Y\}}\!\! n_a(s)\, f^{-}_a(p),
\label{eq:netcharge}
\end{equation}
where $n_a(s)$ counts occurrences of $a$ in $s$ and we use
$\mathrm{p}K_a$ values $K{=}10.5$, $R{=}12.4$, $H{=}6.0$, $D{=}3.9$,
$E{=}4.1$, $C{=}8.3$, $Y{=}10.1$, N-terminus $7.0$, C-terminus $2.34$.
We evaluate $Q(s) := Q(s; 7.4)$.
\end{definition}

\begin{proposition}[Monotonicity of net charge]
\label{prop:charge-monotone}
$Q(s; p)$ is strictly decreasing in $p$ for every nonempty
sequence in this model because its termini are ionizable, with $\lim_{p\to 0} Q = n_K + n_R + n_H + 1$ and
$\lim_{p\to\infty} Q = -(n_D + n_E + n_C + n_Y) - 1$.
\end{proposition}
\begin{proof}
Each term in \eqref{eq:hh} is a logistic function of $p$:
$f^{+}_a(p) = \sigma(\ln 10 \cdot (\mathrm{p}K_a - p))$ is strictly
decreasing in $p$ and $f^{-}_a(p) = \sigma(\ln 10 \cdot (p - \mathrm{p}K_a))$
strictly increasing. $Q(s;p)$ is a positive combination of the decreasing
terms minus a positive combination of the increasing terms, hence strictly decreasing because at least the terminal terms are present. The limits follow
from $\lim_{x\to-\infty}\sigma(x) = 0$ and $\lim_{x\to\infty}\sigma(x) = 1$
applied term by term.
\end{proof}

\begin{definition}[Eisenberg hydrophobic moment]
For window length $w$ and helix angle $\delta$ ($100^\circ$ for an
$\alpha$-helix), the hydrophobic moment of window $s_{i+1}\cdots s_{i+w}$
with Kyte--Doolittle hydropathy values $h_j$ is
\begin{equation}
\mu_i(s; w, \delta) = \frac{1}{w}\left[\left(\sum_{j=1}^{w} h_j \cos j\delta\right)^2
  + \left(\sum_{j=1}^{w} h_j \sin j\delta\right)^2\right]^{1/2},
\label{eq:hmoment}
\end{equation}
and $\mu_H(s) = \max_i \mu_i(s; 11, 100^\circ)$.
\end{definition}

\begin{remark}
Equation \eqref{eq:hmoment} is the length of the mean resultant vector of
the hydropathy values placed at successive residues of a regular helix; it
is maximal exactly when the hydrophobic residues cluster on one face, the
amphipathicity that membrane-active helices exploit. This is the standard
Eisenberg consensus moment.
\end{remark}

\begin{definition}[Pairwise $k$-mer novelty]
For sequence $s$, reference corpus $\mathcal R$, and set $K_k(s)$ of
contiguous $k$-mers, define
\begin{equation}
J_k(s,r)=\frac{|K_k(s)\cap K_k(r)|}{|K_k(s)\cup K_k(r)|},\qquad
\nu_k(s;\mathcal R)=1-\max_{r\in\mathcal R} J_k(s,r).
\label{eq:novelty}
\end{equation}
This is distance from the nearest individual reference under the Jaccard
metric, not the fraction of $k$-mers missing from the union of the corpus.
A high score establishes only weak local sequence similarity; it does not
establish absence from every database or imply a new biological function.
\end{definition}

\begin{proposition}[Nearest-neighbor interpretation]
If $\mathcal R$ is finite and nonempty, $0\leq\nu_k\leq1$,
and $\nu_k(s;\mathcal R)=0$ exactly when at least one reference has the same
set of $k$-mers as $s$.
\end{proposition}
\begin{proof}
$0\leq |A\cap B|\leq |A\cup B|$ for any nonempty finite sets $A,B$,
so $0\leq J_k\leq1$. A finite maximum reaches one precisely if
$K_k(s)=K_k(r)$ for some $r$. Subtracting from one proves both claims.
Notice that equal $k$-mer sets do not imply identical sequences.
\end{proof}

\subsection{Stacked generalization}
\label{sec:stack-math}

Let $b_1, \dots, b_B$ be base scorers trained on the training partition
and let $z(x) = (b_1(x), \dots, b_B(x))^\top \in [0,1]^B$ be the score
vector on the tune partition. The stacking rule is the logistic model
\begin{equation}
g(x) = \sigma(w^\top z(x) + w_0), \qquad
\hat w = \arg\min_{w} \sum_{(x,y) \in \mathcal{T}_{\mathrm{tune}}}
  \ell(y, w^\top z(x) + w_0) + \lambda \|w\|_2^2,
\label{eq:stack}
\end{equation}
with $\ell(y, u) = \log(1 + e^{u}) - yu$.

\begin{proposition}[Convexity and gradient of the stacking objective]
\label{prop:stack-convex}
The objective \eqref{eq:stack} is convex in $(w,w_0)$ and strictly convex
when the augmented tune-partition design matrix $[Z,\mathbf1]$ has full
column rank and all fitted probabilities lie strictly between zero and one, with gradient
$\nabla_w = Z^\top(\sigma(Zw + w_0\mathbf 1) - y) + 2\lambda w$.
\end{proposition}
\begin{proof}
$\ell(y, u)$ is convex in $u$ since
$\partial^2 \ell / \partial u^2 = \sigma(u)(1-\sigma(u)) \ge 0$;
composition with the affine map $(w,w_0) \mapsto w^\top z + w_0$ preserves
convexity, and $\lambda\|w\|_2^2$ adds strict convexity in $w$ for
$\lambda > 0$. Positive logistic curvature plus full rank of the
augmented design makes the Hessian positive definite. The gradient follows from
$\partial \ell / \partial u = \sigma(u) - y$ and the chain rule.
\end{proof}

\subsection{The Matthews correlation coefficient}
\label{sec:mcc-math}

\begin{proposition}[MCC is a Pearson coefficient]
\label{prop:mcc}
For nonconstant binary predictions $\hat y$ and labels $y$ with nonzero
class counts and confusion counts
$TP, TN, FP, FN$,
\begin{equation}
\mathrm{MCC} = \frac{TP \cdot TN - FP \cdot FN}
 {\sqrt{(TP{+}FP)(TP{+}FN)(TN{+}FP)(TN{+}FN)}}
 = \phi(y, \hat y),
\label{eq:mcc}
\end{equation}
the Pearson correlation of the binary vectors; equivalently
$\mathrm{MCC}^2 = \chi^2 / N$ for the $2\times 2$ contingency table.
\end{proposition}
\begin{proof}
For binary vectors, $\mathrm{cov}(y,\hat y) = \frac{TP}{N} -
\frac{(TP{+}FN)(TP{+}FP)}{N^2} = \frac{TP \cdot TN - FP \cdot FN}{N^2}$
after expanding $(TP{+}TN{+}FP{+}FN) = N$, and
$\mathrm{var}(y) = \frac{(TP{+}FN)(TN{+}FP)}{N^2}$,
$\mathrm{var}(\hat y) = \frac{(TP{+}FP)(TN{+}FN)}{N^2}$. Substitution into
$\phi = \mathrm{cov}/\sqrt{\mathrm{var}(y)\mathrm{var}(\hat y)}$ gives
\eqref{eq:mcc}; the $\chi^2$ identity is the standard
$\chi^2 = N\phi^2$ for $2 \times 2$ tables.
\end{proof}

Unlike accuracy, MCC reflects all four confusion-matrix cells; it is
undefined mathematically when a prediction or label vector is constant,
though some software conventions return zero. We report the specified
software implementation for each benchmark comparison.
